\section{微观负二项分支过程与二项稀疏定理}
\label{sec:part02}

\subsection{超传播个体的泊松--伽马混合机制}

在流行病学中，负二项后代分布被广泛用于刻画“超级传播（Superspreading）”现象 \citep{lloydsmith2005superspreading}。本小节从连续型个体传染力的伽马混合出发，显式推导式~\eqref{eq:offspring} 的负二项概率质量函数与理论矩。

\begin{derivationbox}{机制推导 2.1：泊松--伽马混合生成负二项分布}
设感染个体 $i$ 的内在传染力（传代强度）为随机变量 $\nu_i$。由于个体免疫状态、接触网络度数及微环境暴露的异质性，假设 $\nu_i$ 服从均值为 $R_t$、形状参数为 $k$ 的伽马（Gamma）分布：
\begin{equation}
  \nu_i \sim \text{Gamma}\!\left(k,\ \frac{k}{R_t}\right), \qquad f_{\nu}(\nu) = \frac{(k/R_t)^k}{\Gamma(k)} \, \nu^{k-1} e^{-\frac{k}{R_t}\nu}, \quad \nu > 0.
\end{equation}
给定该传染力 $\nu_i$，其产生的继发感染数 $X_i$ 服从均值为 $\nu_i$ 的泊松分布：
\begin{equation}
  X_i \mid \nu_i \sim \text{Poisson}(\nu_i), \qquad P(X_i = x \mid \nu_i) = \frac{e^{-\nu_i} \nu_i^x}{x!}, \quad x = 0, 1, 2, \dots
\end{equation}
对潜变量 $\nu_i$ 积分求边际分布：
\begin{align}
  P(X_i = x \mid R_t) &= \int_0^\infty P(X_i = x \mid \nu) f_\nu(\nu) \, \mathrm{d}\nu \nonumber \\
  &= \int_0^\infty \frac{e^{-\nu} \nu^x}{x!} \frac{(k/R_t)^k}{\Gamma(k)} \nu^{k-1} e^{-\frac{k}{R_t}\nu} \, \mathrm{d}\nu \nonumber \\
  &= \frac{(k/R_t)^k}{x! \, \Gamma(k)} \int_0^\infty \nu^{x+k-1} e^{-\left(1 + \frac{k}{R_t}\right)\nu} \, \mathrm{d}\nu.
\end{align}
引入变量代换 $u = \left(1 + \frac{k}{R_t}\right)\nu = \frac{k+R_t}{R_t}\nu$，利用伽马函数定义 $\int_0^\infty u^{x+k-1} e^{-u} \mathrm{d}u = \Gamma(x+k)$：
\begin{align}
  P(X_i = x \mid R_t) &= \frac{(k/R_t)^k}{x! \, \Gamma(k)} \left( \frac{R_t}{k+R_t} \right)^{x+k} \Gamma(x+k) \nonumber \\
  &= \frac{\Gamma(x+k)}{x! \, \Gamma(k)} \left( \frac{k}{k+R_t} \right)^k \left( \frac{R_t}{k+R_t} \right)^x.
  \label{eq:der-nb-mix}
\end{align}
对照负二项定义式~\eqref{eq:nb-pmf}，式~\eqref{eq:der-nb-mix} 严格等价于参数为 $r = k$ 且成功概率为
\begin{equation}
  p = \frac{k}{k+R_t}, \qquad 1 - p = \frac{R_t}{k+R_t}
\end{equation}
的负二项分布 $\NB\big(k, \frac{k}{k+R_t}\big)$。
\end{derivationbox}

利用全方差公式，个体继发感染数的方差可直接分解为个体传播泊松随机性与环境伽马随机性两部分：
\begin{align}
  \Var(X_i \mid R_t) &= \E\big[\Var(X_i \mid \nu_i)\big] + \Var\big[\E(X_i \mid \nu_i)\big] \nonumber \\
  &= \E(\nu_i) + \Var(\nu_i) \nonumber \\
  &= R_t + \frac{R_t^2}{k}.
  \label{eq:offspring}
\end{align}
显然，当 $k \to \infty$ 时，$\Var(\nu_i) \to 0$，个体传染力退化为常数 $R_t$，后代分布退化为标准泊松分布 $\text{Poisson}(R_t)$。

\subsection{新发总感染数 $I_t$ 的独立求和性质}

当上一代有 $n$ 名互相独立的感染者时，当期新发总感染数为 $I_t = \sum_{i=1}^n X_{i,t}$。

\begin{derivationbox}{推导 2.2：负二项独立和的分布}
由式~\eqref{eq:nb-pgf}，单一个体的后代 PGF 为：
\begin{equation}
  G_X(s) = \left( \frac{p}{1 - (1-p)s} \right)^k, \qquad p = \frac{k}{k+R_t}.
\end{equation}
由于 $X_{1,t}, \dots, X_{n,t}$ 相互独立，和变量 $I_t$ 的概率母函数为各自母函数之积：
\begin{equation}
  G_I(s) = \E\big(s^{\sum_{i=1}^n X_{i,t}}\big) = \prod_{i=1}^n \E(s^{X_{i,t}}) = [G_X(s)]^n = \left( \frac{p}{1 - (1-p)s} \right)^{nk}.
\end{equation}
这正是参数为分散参数 $nk$、成功概率 $p = \frac{k}{k+R_t}$ 的负二项分布。因此，条件在 $n$ 与 $R_t$ 下：
\begin{equation}
  I_t \mid n, R_t \sim \NB\!\left(nk,\ \frac{k}{k+R_t}\right),
\end{equation}
其理论矩为：
\begin{equation}
  \E(I_t \mid n, R_t) = n R_t, \qquad \Var(I_t \mid n, R_t) = n R_t + \frac{n R_t^2}{k}.
\end{equation}
\end{derivationbox}

\subsection{二项稀疏（Binomial Thinning）下负二项族的闭包定理}

在监测系统只能以概率 $\rho$ 报告病例时，观测病例数 $C_t \mid I_t \sim \Bin(I_t, \rho)$。以下证明：二项稀疏操作严格保持负二项分布族的形式不变，仅平移其参数。

\begin{theorembox}{定理 2.1（二项稀疏下的负二项保持定理）}
若总感染数 $I \sim \NB(r, p)$，且观测数 $C \mid I \sim \Bin(I, \rho)$，其中 $\rho \in (0, 1]$，则边际观测数 $C$ 严格服从负二项分布：
\begin{equation}
  C \sim \NB\left(r,\ \frac{p}{p + \rho(1-p)}\right).
\end{equation}
\end{theorembox}

\begin{proof}
利用重期望法则（Law of iterated expectations）与二项分布的概率母函数。已知对于二项分布 $C \mid I \sim \Bin(I, \rho)$，其条件母函数为：
\begin{equation}
  \E(s^C \mid I) = \sum_{c=0}^I \binom{I}{c} (\rho s)^c (1-\rho)^{I-c} = (1 - \rho + \rho s)^I.
\end{equation}
因此，$C$ 的边际概率母函数 $G_C(s)$ 为：
\begin{equation}
  G_C(s) = \E\big[\E(s^C \mid I)\big] = \E\big[(1 - \rho + \rho s)^I\big] = G_I(1 - \rho + \rho s).
  \label{eq:thin-compound}
\end{equation}
将 $I \sim \NB(r, p)$ 的母函数表达式代入式~\eqref{eq:thin-compound}，将变量 $s$ 整体替换为 $1 - \rho + \rho s$：
\begin{align}
  G_C(s) &= \left( \frac{p}{1 - (1-p)(1 - \rho + \rho s)} \right)^r \nonumber \\
  &= \left( \frac{p}{1 - (1-p)(1-\rho) - (1-p)\rho s} \right)^r.
  \label{eq:thin-expand}
\end{align}
定义辅助分母常数项：
\begin{equation}
  q \equiv 1 - (1-p)(1-\rho) = 1 - (1 - \rho - p + p\rho) = p + \rho(1-p).
\end{equation}
在式~\eqref{eq:thin-expand} 的分母中提出因子 $q$：
\begin{equation}
  G_C(s) = \left( \frac{p}{q - (1-p)\rho s} \right)^r = \left( \frac{p/q}{1 - \frac{(1-p)\rho}{q} s} \right)^r.
  \label{eq:thin-factored}
\end{equation}
注意到：
\begin{equation}
  1 - \frac{p}{q} = \frac{q - p}{q} = \frac{p + \rho(1-p) - p}{q} = \frac{(1-p)\rho}{q}.
\end{equation}
将此关系代入式~\eqref{eq:thin-factored}，得：
\begin{equation}
  G_C(s) = \left( \frac{p'}{1 - (1 - p')s} \right)^r, \qquad \text{其中 } p' \equiv \frac{p}{q} = \frac{p}{p + \rho(1-p)}.
\end{equation}
对照负二项分布概率母函数的定义式~\eqref{eq:nb-pgf}，母函数形式完全相同，分散参数仍为 $r$，成功概率由 $p$ 变换为 $p'$。由于概率母函数与离散概率分布之间存在一一对应性，定理得证。
\end{proof}

\subsection{报告病例数 $C_t$ 的精确条件矩推导}

\begin{derivationbox}{推导 2.3：代入流行病学参数}
在我们的模型中，分散参数为 $r = nk$，原成功概率为 $p = \frac{k}{k + R_t}$，故 $1 - p = \frac{R_t}{k + R_t}$。计算稀疏后的分母常数 $q$ 与新参数 $p'$：
\begin{equation}
  q = p + \rho(1-p) = \frac{k}{k+R_t} + \frac{\rho R_t}{k+R_t} = \frac{k + \rho R_t}{k + R_t},
\end{equation}
\begin{equation}
  p' = \frac{p}{q} = \frac{\frac{k}{k+R_t}}{\frac{k+\rho R_t}{k+R_t}} = \frac{k}{k + \rho R_t}.
\end{equation}
由此得出核心结论：
\begin{equation}
  \boxed{\, C_t \mid n, R_t \sim \NB\left(nk,\ \frac{k}{k + \rho R_t}\right) \,}
  \label{eq:der-ct-dist}
\end{equation}
定义 $m \equiv \rho n$ 为预期报告规模。应用负二项分布矩公式式~\eqref{eq:nb-mean} 与式~\eqref{eq:nb-var}：
\begin{align}
  \E(C_t \mid n, R_t) &= \frac{r(1-p')}{p'} = nk \cdot \frac{\frac{\rho R_t}{k+\rho R_t}}{\frac{k}{k+\rho R_t}} = \rho n R_t = m R_t, \\
  \Var(C_t \mid n, R_t) &= \frac{r(1-p')}{(p')^2} = nk \cdot \frac{\rho R_t}{k+\rho R_t} \cdot \left(\frac{k+\rho R_t}{k}\right)^2 \nonumber \\
  &= \rho n R_t \left(1 + \frac{\rho R_t}{k}\right) = \rho n R_t + \frac{\rho^2 n R_t^2}{k} \nonumber \\
  &= m R_t + \frac{\rho m R_t^2}{k}.
\end{align}
进而得到相对方差（平方变异系数 $\CV^2$）：
\begin{equation}
  \CV^2(C_t \mid n, R_t) = \frac{\Var(C_t \mid n, R_t)}{[\E(C_t \mid n, R_t)]^2} = \frac{\rho n R_t + \rho^2 n R_t^2 / k}{(\rho n R_t)^2} = \frac{1}{\rho n R_t} + \frac{1}{nk} = \frac{1}{m R_t} + \frac{1}{nk}.
  \label{eq:der-cv2}
\end{equation}
\end{derivationbox}

\begin{insightbox}{物理与流行病学洞见 2.1}
式~\eqref{eq:der-cv2} 清晰揭示了观测噪声的内在结构：
\begin{enumerate}
  \item 第一项 $\frac{1}{\rho n R_t} = \frac{1}{m R_t}$ 为\textbf{泊松/二项抽样噪声}。提高报告率 $\rho$ 或增大祖代 $n$ 都能将其压缩；
  \item 第二项 $\frac{1}{nk}$ 为\textbf{超传播个体方差}。请注意：该项分母中\textbf{完全不包含报告率 $\rho$}！这意味着：即使建立 100\% 完美的报告系统（$\rho = 1$），由于群体内超级传播者的存在，每代传播的相对方差依然保留底板 $\frac{1}{nk}$。要消除此项，唯一途径是扩大流行被追踪祖代的绝对人数 $n$。
\end{enumerate}
\end{insightbox}
